% ===================================================================
%  TABEL Nr. 4 — Volwassen leeftijd en ouderdom.
%  Traduction 2026 d'apres texto/io, controlee sur texto/fr et
%  texto/es. Consigne : voir texto/nl/10-tabel-01.tex.
% ===================================================================

%%K t04-apar-1 apar
\VUsaut{4.0mm}
\VUtitre{15.0pt}{60}{TABEL Nr. 4}
\VUsaut{1.6mm}
\VUtitre{11.4pt}{0}{Volwassen leeftijd en ouderdom.}
\VUsaut{3.2mm}

%%K t04-tit-1 sub
\VUcentre{11.4pt}{0}{\textit{Eerste tafereel.}}
\VUsaut{1.6mm}
\VUtitre{11.4pt}{0}{Het bruiloftsmaal.}
\VUsaut{2.6mm}

%%K t04-tit-2 sub
\VUpk{10.2pt}{40}{De herfst.}
\VUsaut{3.0mm}

%%K t04-c1-01-1 p
1. --- De jongelui zijn groter geworden. Een van hen, Jan, trouwt. Wij wonen het
\VUgras{bruiloftsmaal} bij dat de families van de echtgenoten om dezelfde tafel
verenigt.

%%K t04-c1-02-1 p
2. --- Wij zijn in de \VUgras{herfst}, in de \VUgras{maand september}. De koude wind van
\VUgras{oktober} en van \VUgras{november} heeft het land nog niet van zijn groene loof
beroofd. De avondlucht is lauw en geurig; daarom vindt het diner plaats onder de
\VUgras{veranda} \textsuperscript{(8)} bij de tuin.

%%K t04-c1-03-1 p
3. --- In de verte, op de \VUgras{heuvel} \textsuperscript{(3)}, ligt het huis van de ouders
van Jan, een sierlijk \VUgras{kasteel} \textsuperscript{(1)} verborgen in het groen; mooie
wijngaarden, die een uitstekende wijn opleveren, bedekken de helling van de heuvel. De
wijnbouwer bewerkt en verzorgt ze.

%%K t04-c1-04-1 p
4. --- Boven de wijngaarden vliegt een vlucht \VUgras{patrijzen} \textsuperscript{(2)} die
verschrikt zijn door het naderen van de
\VUgras{jachtopziener} \textsuperscript{(5)}. Deze doorzoekt, met zijn
\VUgras{geweer} onder de arm en een \VUgras{weitas over de schouder}, het
\VUgras{struikgewas} \textsuperscript{(4)} met zijn \VUgras{hond}
\textsuperscript{(6)}. Hij houdt toezicht op het landgoed en belet de stropers het wild uit te
roeien.

%%K t04-c1-05-1 p
5. --- Buiten de \VUgras{veranda} klimt een \VUgras{leiwingerd} waaraan dichte
\VUgras{druiventrossen} \textsuperscript{(7)} hangen.

%%K t04-c1-06-1 p
6. --- De mooie herfstbloemen, die de \VUgras{tafel} \textsuperscript{(16)} versieren, zijn
\VUgras{chrysanten} \textsuperscript{(9)}; de \VUgras{vader} \textsuperscript{(10)} van de
bruid heeft er een in het knoopsgat van zijn rok gestoken.

%%K t04-c1-07-1 p
7. --- De maaltijd loopt ten einde. De \VUgras{bediende} \textsuperscript{(11)} begint de
nagerechtschotels te brengen en zal weldra de verschillende delen van het couvert wegnemen,
\VUgras{lepels} \textsuperscript{(14)}, \VUgras{vorken} \textsuperscript{(12)},
\VUgras{messen} \textsuperscript{(13)},
\VUgras{schalen} \textsuperscript{(15)} die men niet meer nodig heeft.

%%K t04-tit-3 sub
\VUpk{10.2pt}{40}{De verwantschap.}
\VUsaut{3.0mm}

%%K t04-c1-08-1 p
8. --- Allen zien er vrolijk uit, en de glimlach waarmee de \VUgras{moeder}
\textsuperscript{(17)} van de bruidegom haar kinderen aankijkt bewijst haar geluk.

%%K t04-c1-09-1 p
9. --- Dat huwelijk verenigt twee rijke families van de streek. Jan, de
\VUgras{echtgenoot} \textsuperscript{(18)}, is de \VUgras{zoon} van een groot
grondbezitter en hij wordt de \VUgras{schoonzoon} van een bekwaam industrieel.

%%K t04-c1-10-1 p
10. --- De \VUgras{echtgenoten} zijn jong en toch waren zij al lang
\VUgras{verloofd}. De \VUgras{jonge vrouw} \textsuperscript{(19)}, Martha, is
bekoorlijk in haar \VUgras{witte jurk} versierd met oranjebloesem.

%%K t04-c1-11-1 p
11. --- Haar \VUgras{schoonvader} \textsuperscript{(20)} is heel gelukkig zo'n lieve
\VUgras{schoondochter} te hebben, en de \VUgras{schoonmoeder} \textsuperscript{(21)} van Jan
glimlacht als zij haar \VUgras{dochter} zo mooi ziet.

%%K t04-c1-12-1 p
12. --- Aan het uiteinde van de tafel zitten de overige \VUgras{genodigden}
\textsuperscript{(22)}, \VUgras{bloedverwanten, vrienden, neven, nichten}. Een
\VUgras{gerant} \textsuperscript{(23)}, met een servet onder de arm, bedient hen
ijverig.

%%K t04-tit-4 sub
\VUpk{10.2pt}{40}{De vrienden.}
\VUsaut{3.0mm}

%%K t04-c1-13-1 p
13. --- Aan de rechterkant van de tafel zit een \VUgras{jongedame} \textsuperscript{(24)},
\VUgras{vriendin} van de jonge bruid, daarna de
\VUgras{broer} van deze laatste, die haar \VUgras{bruidsjonker}
\textsuperscript{(25)} is. Hij praat met zijn \VUgras{bruidsmeisje} \textsuperscript{(26)}
zuster van de bruidegom en die door dat huwelijk de
\VUgras{schoonzuster} van Martha wordt. Zij is een bekoorlijk meisje. Zij heeft
mooie \VUgras{juwelen}, een \VUgras{parelsnoer}, en een gouden
\VUgras{armband}.

%%K t04-c1-14-1 p
14. --- Naast hen is de heer die opstaat en zijn glas heft een \VUgras{vriend}
\textsuperscript{(27)} van het gezin. Hij is een van de \VUgras{getuigen van de bruidegom}.
Hij \VUgras{brengt een heildronk uit}, drinkt op de
\VUgras{gezondheid} van de jonge echtgenoten en wenst hun een lang geluk. Hij
draagt een gelegenheidskostuum, een rok met \VUgras{lakschoenen} \textsuperscript{(28)}. Hij
begeleidt de \VUgras{grootmoeder} \textsuperscript{(29)} die \VUgras{weduwe} is.

%%K t04-c1-15-1 p
15. --- De jongeman in \VUgras{rok} \textsuperscript{(31)}, met een dun zwart snorretje, is de
\VUgras{zwager} \textsuperscript{(30)} van Jan. Hij is een vooraanstaand ingenieur. Hij zit
naast een heel sierlijke \VUgras{jonge dame} \textsuperscript{(33)}, de \VUgras{oudere zuster}
van Martha en de moeder van het lieve meisje dat, aan de arm van haar neef, een bloem komt
aanbieden en hem
\VUgras{goedenavond} komt \VUgras{wensen}. Die dame heeft heel mooie
\VUgras{oorbellen} en een sierlijk \VUgras{kapsel}.

%%K t04-c1-16-1 p
16. --- Het neefje, zo vreemd in zijn \VUgras{kiel} \textsuperscript{(36)} en zijn wijde
\VUgras{broek} \textsuperscript{(37)} is heel beklagenswaardig. Hij is een \VUgras{wees}
\textsuperscript{(35)}. Heel jong verloor hij zijn vader en zijn moeder. Zijn oom is zijn
\VUgras{voogd} \textsuperscript{(38)}, en is vrijgezel gebleven om het arme kind op te voeden.
Hij is een goedhartig man. Hij is een beetje kaal en heel bijziend; hij gebruikt een
\VUgras{knijpbril}.

%%K t04-tit-5 sub
\VUsaut{2.4mm}
\VUpk{10.2pt}{40}{Het nagerecht.}
\VUsaut{3.0mm}

%%K t04-c1-17-1 p
17. --- Achter de gasten, op een tafel bedekt met een \VUgras{tafellaken}, is het
\VUgras{nagerecht} \textsuperscript{(39)} klaargezet. In een grote schaal ligt fruit:
\VUgras{perziken} \textsuperscript{(40)}, \VUgras{peren} \textsuperscript{(41)},
\VUgras{appels} \textsuperscript{(42)},
\VUgras{abrikozen} \textsuperscript{(43)} en \VUgras{druiven}
\textsuperscript{(44)}. Dichtbij, \VUgras{ananassen} \textsuperscript{(45)}, een mooie
\VUgras{taart} \textsuperscript{(46)},
\VUgras{likeurflessen} \textsuperscript{(48)} en \VUgras{champagne}
\textsuperscript{(49)} en ook stapels schone \VUgras{borden} \textsuperscript{(47)}.

%%K t04-c1-18-1 p
18. --- De \VUgras{kelnermeester} \textsuperscript{(50)} ontkurkt een
\VUgras{fles} \textsuperscript{(52)} oude wijn met een \VUgras{kurkentrekker}
\textsuperscript{(51)}. De \VUgras{kurk} \textsuperscript{(53)} lijkt erg moeilijk uit te
trekken. Die kelnermeester heeft een \VUgras{servet} \textsuperscript{(54)} onder de arm.

%%K t04-c1-19-1 p
19. --- Na het diner zullen de heren naar de rookkamer gaan, en de dames zullen zich voor het
bal gaan klaarmaken.

%%K t04-tit-6 sub
\VUsaut{5.0mm}
\VUcentre{11.4pt}{0}{\textit{Tweede tafereel.}}
\VUsaut{1.6mm}
\VUpk{11.4pt}{40}{De verjaardag van grootvader.}
\VUsaut{2.6mm}

%%K t04-tit-7 sub
\VUpk{10.2pt}{40}{Het bezoek.}
\VUsaut{3.0mm}

%%K t04-c2-01-1 p
1. --- Tien jaar zijn voorbijgegaan, de ouders van Jan zijn oud geworden en houden veel van
hun drie kleinkinderen, twee \VUgras{kleinzoons} \textsuperscript{(2)} en een
\VUgras{kleindochter} \textsuperscript{(3)}. Omdat vandaag
\VUgras{de verjaardag van grootvader} is, komen de kinderen hun een gelukkige naamdag wensen.

%%K t04-c2-02-1 p
2. --- De kleine Piet is nog heel jong, hij is pas drie jaar. Hij ziet de
\VUgras{kat} \textsuperscript{(1)} naderen. Hij wil haar mooie zijdeachtige
\VUgras{vacht} en haar lange \VUgras{staart} strelen, hij bedenkt niet, dat zich
onder de sierlijke \VUgras{poten} boosaardige spitse klauwen verbergen die hem misschien
zullen krabben.

%%K t04-c2-03-1 p
3. --- Zijn zusje Gabriëlle, dat hem de hand geeft, is veel ernstiger en Marcel met zijn grote
\VUgras{mantel} \textsuperscript{(4)} en zijn leren
\VUgras{riem} \textsuperscript{(5)} ziet er wat mannelijker uit, ondanks zijn
korte broek die zijn \VUgras{kousen} \textsuperscript{(6)} laat zien.

%%K t04-c2-04-1 p
4. --- Hun vader heeft zijn dikke \VUgras{winterjas} \textsuperscript{(7)} met
\VUgras{bontkraag} \textsuperscript{(8)} aangetrokken en in zijn \VUgras{zak}
\textsuperscript{(9)} heeft hij een sjaal. Zijn vrouw heeft haar vilten hoed met
\VUgras{veren} \textsuperscript{(10)} op, haar \VUgras{jasje}
\textsuperscript{(11)}, haar \VUgras{bontboa} \textsuperscript{(12)} en haar
\VUgras{mof} \textsuperscript{(13)}.

%%K t04-c2-05-1 p
5. --- Het is erg koud, want het is winter. De \VUgras{sneeuw} \textsuperscript{(19)} is
de hele dag gevallen en de daken van de huizen zijn helemaal wit. Tegelijk met de
\VUgras{nacht} wordt de koude nog scherper. Aan de hemel schijnen de \VUgras{maan}
\textsuperscript{(17)} en de \VUgras{sterren} \textsuperscript{(18)} en doorbreken
\VUgras{de duisternis}.

%%K t04-tit-8 sub
\VUsaut{4.2mm}
\VUpk{10.2pt}{40}{De ziekte.}
\VUsaut{3.0mm}

%%K t04-c2-06-1 p
6. --- De arme \VUgras{blinde} \textsuperscript{(20)} die het plein oversteekt voor
\VUgras{de apotheek} \textsuperscript{(16)} geleid door zijn
\VUgras{poedel} \textsuperscript{(21)}, is erg ongelukkig. Justine, de
\VUgras{dienstbode} \textsuperscript{(14)}, brengt uit de keuken een
\VUgras{kom} \textsuperscript{(15)} erg warme \VUgras{kruidenthee} die rijkelijk
gesuikerd is.

%%K t04-c2-07-1 p
7. --- De \VUgras{grootmoeder} \textsuperscript{(23)} die haar
\VUgras{lorgnet} \textsuperscript{(26)} op de tafel wil zoeken, om het
\VUgras{recept} \textsuperscript{(27)} te lezen dat de \VUgras{dokter}
\textsuperscript{(25)} zojuist heeft geschreven, staat op uit haar leunstoel en steunt daarbij
op haar \VUgras{stok} \textsuperscript{(24)}.

%%K t04-c2-08-1 p
8. --- Dikwijls lijdt zij aan \VUgras{ongesteldheden, duizelingen, verkoudheden, kiespijn},
haar benen zijn zwak en haar ogen zijn niet meer zo goed. Toch spot zij graag met haar oude
\VUgras{dokter} die haar altijd een
\VUgras{drankje} \textsuperscript{(28)} wil voorschrijven. Vandaag is zij heel
vrolijk omdat zij haar jonge familie om zich heen heeft.

%%K t04-c2-09-1 p
9. --- Het is behaaglijk in de goed gesloten kamer. In de marmeren
\VUgras{schoorsteen} \textsuperscript{(29)} vlamt een \VUgras{prachtig vuur}
\textsuperscript{(30)}, en men voelt zich gelukkig in het zachte \VUgras{licht} van de
\VUgras{lamp} \textsuperscript{(31)} die men zojuist heeft aangestoken.

%%K t04-tit-9 sub
\VUsaut{4.2mm}
\VUpk{10.2pt}{40}{De grootvader.}
\VUsaut{3.0mm}

%%K t04-c2-10-1 p
10. --- Zelfs de \VUgras{grootvader} \textsuperscript{(33)} is vergeten, dat hij ziek was, dat
zijn \VUgras{reumatiek} hem aan zijn \VUgras{leunstoel} \textsuperscript{(34)} kluistert en
dat hij gisteren nog \VUgras{koorts en flauwtes} had. Hij herinnert zich niet meer, dat hij de
afgelopen winter aan
\VUgras{longontsteking} leed en dat een schorre en pijnlijke \VUgras{hoest} hem
heel deed lijden.

%%K t04-c2-10-2 p
En toch werd hij erg moeilijk genezen. Voortaan gaat het hem wat beter. Maar zijn been dat
hij op een \VUgras{kussen} \textsuperscript{(38)} legt dat op een klein \VUgras{krukje}
\textsuperscript{(39)} staat doet hem ook pijn.

%%K t04-c2-11-1 p
11. --- Wanneer hij zorgvuldig gewikkeld is in zijn warme \VUgras{kamerjas}
\textsuperscript{(35)} met dikke paarlemoeren \VUgras{knopen} \textsuperscript{(36)} en
wanneer hij naast zich, om hem de krant voor te lezen, het kleine \VUgras{weesmeisje}
\textsuperscript{(40)} heeft dat een witte
\VUgras{muts}, een blauwe \VUgras{jurk} \textsuperscript{(42)} en een
\VUgras{pelerine} \textsuperscript{(41)} draagt, klaagt hij niet.

%%K t04-c2-12-1 p
12. --- Elk jaar, wanneer de \VUgras{kalender} \textsuperscript{(43)} opnieuw de
\VUgras{datum} van de eerste \VUgras{december} aanwijst, wacht hij ongeduldig op
zijn \VUgras{kleinkinderen}, want hij weet, dat zij die belangrijke
\VUgras{datum} niet zullen vergeten.

%%K t04-c2-13-1 p
13. --- Hij herinnert zich met ontroering de heel lang vervlogen tijd toen hij persoonlijk de
roemrijke keizerlijke \VUgras{maarschalk} een gelukkige naamdag kwam wensen, wiens
\VUgras{portret} \textsuperscript{(44)} aan de muur hangt, en die de \VUgras{overgrootvader}
van zijn kinderen is.
\VUsaut{6.0mm}
\VUfilet{22mm}
